%% 6_discussion
\section{Limitations}
\label{sec:limitations}

We state the boundaries of these claims plainly, because several of them constrain
the conclusion.

\noindent\textbf{One simulated corpus.} The plateau in \cref{fig:plateau} is
measured on RadioMapSeer, which is ray traced with a single propagation engine. A
plateau shared by nine architectures on one corpus is strong evidence about that
corpus and suggestive about the problem. It would be stronger across two
independently generated corpora, and we have not shown that.

\noindent\textbf{Two real sites, not four.} The corpus spans four buildings but
only two yield cells with enough co-located repeats to evaluate on. The wall effect
and its control therefore rest on one treatment site and one control site. A second
walled site would turn a clean contrast into a replicated one.

\noindent\textbf{The real evaluation cannot test extrapolation.} Reference points
are clustered, so holding one out asks for interpolation among near neighbours.
This is the honest limit of measurements collected without a dense ground-truth
survey, and it is why we do not read the classical win in \cref{tab:real} as
evidence against learned reconstruction in general. Settling that question needs a
densely surveyed floor, which we did not have.

\noindent\textbf{Uncertainty does not transfer.} Monte Carlo dropout is reasonably
ordered on simulated data, with an error-uncertainty correlation of $0.55$, but on
real measurements the nominal one-sigma interval covers $16.1\%$ of held-out points
where a calibrated Gaussian would cover $68\%$. Recalibration restores coverage
only by widening intervals roughly elevenfold, which is a change in sharpness large
enough that we do not report the recalibrated model as usable.

\noindent\textbf{A negative result is not an impossibility proof.}
\Cref{sec:txfree_result} shows one transmitter-free construction fails and explains
why. It does not show that no such construction exists, and \cref{sec:conclusion}
argues the opposite is worth attempting.
